\begin{figure}[htbp]
\centering
\begin{tikzpicture}[scale=1.1]
    % Define the s-plane extent
    \def\xmin{-5}
    \def\xmax{3}
    \def\ymin{-3.5}
    \def\ymax{3.5}

    % Shade the stable region (LHP)
    \fill[UofSCHorseshoe, opacity=0.15] (\xmin,\ymin) rectangle (0,\ymax);

    % Shade the unstable region (RHP)
    \fill[UofSCRose, opacity=0.10] (0,\ymin) rectangle (\xmax,\ymax);

    % Grid
    \draw[help lines, gray!30] (\xmin,\ymin) grid (\xmax,\ymax);

    % Axes
    \draw[thick, ->] (\xmin,0) -- (\xmax+0.3,0) node[right, font=\small] {$\Real(s) = \sigma$};
    \draw[thick, ->] (0,\ymin) -- (0,\ymax+0.3) node[above, font=\small] {$\Imag(s) = \jj\omega$};

    % Imaginary axis (stability boundary) - emphasized
    \draw[very thick, UofSCCongaree] (0,\ymin) -- (0,\ymax);

    % Region labels
    \node[font=\bfseries, UofSCHorseshoe] at (-2.5, 3.0) {STABLE};
    \node[font=\small, UofSCHorseshoe] at (-2.5, 2.5) {(Left Half-Plane)};
    \node[font=\bfseries, Rose] at (1.5, 3.0) {UNSTABLE};
    \node[font=\small, Rose] at (1.5, 2.5) {(Right Half-Plane)};

    % Stability boundary label
    \node[font=\footnotesize, UofSCCongaree, rotate=90, anchor=south] at (0.25, 0) {Stability Boundary};

    % ========== Example Poles ==========

    % Stable real pole (overdamped component)
    \node[circle, fill, UofSCGarnet, minimum size=8pt] at (-3, 0) {};
    \node[below, font=\footnotesize, UofSCGarnet] at (-3, -0.3) {$p_1 = -3$};
    \node[above right, font=\tiny, align=left] at (-2.8, 0.1) {Stable\\(exponential decay)};

    % Stable complex conjugate poles (underdamped)
    \node[circle, fill, UofSCAtlantic, minimum size=8pt] at (-1.5, 2) {};
    \node[circle, fill, UofSCAtlantic, minimum size=8pt] at (-1.5, -2) {};
    \node[right, font=\footnotesize, UofSCAtlantic] at (-1.4, 2.2) {$p_2 = -1.5 + 2\jj$};
    \node[right, font=\footnotesize, UofSCAtlantic] at (-1.4, -2.2) {$p_2^* = -1.5 - 2\jj$};
    \node[left, font=\tiny, align=right] at (-1.6, 1.5) {Stable\\(damped oscillation)};

    % Marginally stable poles (on imaginary axis)
    \node[circle, fill, UofSCCongaree, minimum size=8pt] at (0, 1) {};
    \node[circle, fill, UofSCCongaree, minimum size=8pt] at (0, -1) {};
    \node[left, font=\footnotesize, UofSCCongaree] at (-0.15, 1.2) {$p_3 = \jj$};
    \node[left, font=\footnotesize, UofSCCongaree] at (-0.15, -1.2) {$p_3^* = -\jj$};
    \node[right, font=\tiny, align=left] at (0.15, 0.5) {Marginally stable\\(sustained oscillation)};

    % Unstable real pole
    \node[circle, fill, Rose, minimum size=8pt] at (1, 0) {};
    \node[above, font=\footnotesize, Rose] at (1, 0.3) {$p_4 = +1$};
    \node[below, font=\tiny, align=center] at (1, -0.3) {Unstable\\(exponential growth)};

    % Unstable complex conjugate poles
    \node[circle, fill, UofSCDarkGarnet, minimum size=8pt] at (0.5, 2.5) {};
    \node[circle, fill, UofSCDarkGarnet, minimum size=8pt] at (0.5, -2.5) {};
    \node[right, font=\footnotesize, UofSCDarkGarnet] at (0.6, 2.7) {$p_5 = 0.5 + 2.5\jj$};
    \node[right, font=\footnotesize, UofSCDarkGarnet] at (0.6, -2.7) {$p_5^* = 0.5 - 2.5\jj$};
    \node[left, font=\tiny, align=right] at (0.4, 2.0) {Unstable\\(growing oscillation)};

    % ========== Annotations ==========

    % Damping annotation
    \draw[UofSCAtlantic, thick, <->] (-1.5, 0.1) -- (0, 0.1);
    \node[above, font=\tiny, UofSCAtlantic] at (-0.75, 0.15) {$|\sigma|$ = decay rate};

    % Frequency annotation
    \draw[UofSCAtlantic, thick, <->] (-1.4, 0) -- (-1.4, 2);
    \node[right, font=\tiny, UofSCAtlantic] at (-1.35, 1) {$\omega_d$};

    % Natural frequency annotation
    \draw[UofSC70Black, thin, dashed] (0, 0) -- (-1.5, 2);
    \node[font=\tiny, UofSC70Black, rotate=53] at (-0.9, 1.3) {$\omega_n$};

    % Damping ratio angle
    \draw[UofSCHoneycomb, thick] (-0.4, 0) arc (180:127:0.4);
    \node[font=\tiny, UofSCHoneycomb] at (-0.6, 0.25) {$\theta$};

    % Legend box
    \begin{scope}[xshift=-4.5cm, yshift=-2.8cm]
        \draw[thick, UofSC70Black] (0,0) rectangle (2.8,1.6);
        \node[font=\footnotesize\bfseries] at (1.4, 1.4) {Pole Types};
        \node[circle, fill, UofSCGarnet, minimum size=5pt] at (0.2, 1.0) {};
        \node[right, font=\tiny] at (0.35, 1.0) {Real (stable)};
        \node[circle, fill, UofSCAtlantic, minimum size=5pt] at (0.2, 0.7) {};
        \node[right, font=\tiny] at (0.35, 0.7) {Complex (stable)};
        \node[circle, fill, UofSCCongaree, minimum size=5pt] at (0.2, 0.4) {};
        \node[right, font=\tiny] at (0.35, 0.4) {Imaginary (marginal)};
        \node[circle, fill, Rose, minimum size=5pt] at (1.5, 1.0) {};
        \node[right, font=\tiny] at (1.65, 1.0) {Real (unstable)};
        \node[circle, fill, UofSCDarkGarnet, minimum size=5pt] at (1.5, 0.7) {};
        \node[right, font=\tiny] at (1.65, 0.7) {Complex (unstable)};
    \end{scope}

\end{tikzpicture}
\caption{Stability regions in the complex $s$-plane. The left half-plane (LHP, shaded green) contains stable poles where $\Real(s) < 0$, producing decaying responses. The right half-plane (RHP, shaded pink) contains unstable poles where $\Real(s) > 0$, producing unbounded responses. The imaginary axis (Congaree, emphasized) is the stability boundary: poles on this axis yield marginally stable systems with sustained oscillations. Example poles demonstrate different behaviors: real stable pole ($p_1$, exponential decay), complex stable pair ($p_2, p_2^*$, damped oscillation), imaginary pair ($p_3, p_3^*$, sustained oscillation), real unstable pole ($p_4$, exponential growth), and complex unstable pair ($p_5, p_5^*$, growing oscillation).}
\label{fig:stability_regions}
\end{figure}
